% ECCV 2026 Workshop X-Reason submission
% Download the ECCV 2026 style files from https://eccv.ecva.net/Conferences/2026
% and place eccv2026.sty + splncs04.bst in this directory.
%
% To compile:  pdflatex main && bibtex main && pdflatex main && pdflatex main

\documentclass[runningheads]{llncs}

\usepackage{eccv2026}           % comment out if style file not yet downloaded
\usepackage{graphicx}
\usepackage{amsmath,amssymb}
\usepackage{booktabs}
\usepackage{multirow}
\usepackage{xcolor}
\usepackage{hyperref}
\usepackage{microtype}

% Convenience macros for placeholders
\newcommand{\TBD}[1]{\textcolor{red}{\textbf{[TBD: #1]}}}
\newcommand{\TODO}[1]{\textcolor{blue}{\textbf{[TODO: #1]}}}
\newcommand{\numimages}{200}      % update when dataset is fixed
\newcommand{\numactions}{6}
\newcommand{\numllms}{3}

\begin{document}

\title{Can Frontier Vision-Language Models Reason About the Interactable World?\\
A Structured Affordance Benchmark}

\titlerunning{VLM Affordance Reasoning Benchmark}

\author{
  \TBD{Author Names}
}
\authorrunning{\TBD{Author et al.}}
\institute{\TBD{Institution}}

\maketitle

% ============================================================
\begin{abstract}
Understanding what actions objects support—their \emph{affordances}—is a
prerequisite for intelligent interaction with the physical world. We ask
whether frontier Vision-Language Models (VLMs) have internalised this
capacity: not just \emph{what} an object is, but \emph{what can be done
with it}, \emph{whether} that is appropriate, and \emph{why}. We
introduce \textbf{AffBench}, a benchmark that evaluates three frontier
VLMs—GPT-4.1, Claude~Sonnet~4.6, and Gemini~2.5~Pro—on structured
affordance reasoning across \numimages{} real-world scenes and
\numactions{} action categories. Our pipeline uses the Segment Anything
Model (SAM)~\cite{kirillov2023sam} for annotation-free instance discovery
and queries each model under a novel \textbf{7-way affordance taxonomy}
spanning physical, social, and safety dimensions: \emph{Positive},
\emph{Firmly Negative}, \emph{Object Non-functional}, \emph{Physical
Obstacle}, \emph{Socially Awkward}, \emph{Socially Forbidden}, and
\emph{Dangerous}. Exception categories additionally require grounded
one-sentence explanations and consequences, enabling text-quality
evaluation. We further compare area-ranked instance selection against a
saliency-guided variant using OOAL affordance maps~\cite{ooal}, and
report inter-model agreement as a ground-truth-free reliability signal.
Our results show that VLMs handle coarse affordance distinctions well
but diverge on fine-grained exception attribution—the harder and more
practically relevant judgement. Code and data are publicly available.
\end{abstract}

\keywords{Affordance reasoning \and Vision-language models \and
Scene understanding \and Segment Anything \and Interactable world}


% ============================================================
\section{Introduction}
\label{sec:intro}

Gibson~\cite{gibson1979ecological} described \emph{affordances} as the
action possibilities that an environment offers to an agent: a chair
\emph{affords} sitting, a knife \emph{affords} cutting, and a crowded
dining table may \emph{forbid} lying on for social reasons. This concept,
grounded in the coupling between perception and action, lies at the heart
of embodied intelligence. An agent that can identify objects but not
reason about what it can \emph{do} with them is fundamentally limited in
any interactive setting—from household robotics and assistive AI to
augmented reality and autonomous planning.

The emergence of powerful Vision-Language Models (VLMs) raises a
compelling question for the field: do these models understand the
\emph{interactable world}? Trained on internet-scale image-text
corpora~\cite{gpt4,claude,gemini}, VLMs encode encyclopaedic semantic
knowledge. They identify objects, describe scenes, and answer factual
questions with impressive fluency. But structured affordance reasoning is
qualitatively harder: it requires simultaneously integrating visual
appearance, physical plausibility, object functionality, social norms, and
safety constraints—and attributing the \emph{reason} an action is
appropriate or forbidden. A VLM that labels a broken chair as ``chair''
may still fail to recognise that sitting on it is dangerous; a model that
knows a religious artefact by name may not understand that certain actions
towards it are socially forbidden.

Prior work on VLMs and affordances has largely focused on the
\emph{where}: localising interaction regions on objects given an action
query~\cite{affordancellm,luo2022grounded}. A smaller body of work
addresses binary \emph{can/cannot} affordance
judgements~\cite{gao2023physically,context_aff2026}. The richer question
of \emph{which type of constraint} applies—physical, social, or safety—has
received almost no systematic attention. This distinction is practically
important: a robot that refuses to pick up a glass because it
``cannot'' (code~1: Firmly Negative) behaves very differently from one
that refuses because the glass is cracked (code~2: Non-functional) or
because picking it up would be culturally inappropriate in context
(code~5: Socially Forbidden). Getting the \emph{reason} right, not just
the binary outcome, is essential for interpretable, safe, and contextually
appropriate AI behaviour.

\noindent\textbf{This work.} We present \textbf{AffBench}, a benchmark
for evaluating frontier VLMs on structured affordance reasoning in
real-world scenes. Our pipeline requires no manual instance annotation: we
use the Segment Anything Model (SAM)~\cite{kirillov2023sam} to propose
instances, optionally rank them by action-specific saliency maps from
OOAL~\cite{ooal}, and query each model with a structured prompt that
enforces our 7-way taxonomy. We evaluate GPT-4.1, Claude~Sonnet~4.6, and
Gemini~2.5~Pro across \numactions{} action categories and \numimages{}
scenes, using inter-model agreement as a GT-free reliability signal in
addition to ground-truth comparison against ADE-Affordance
labels~\cite{ade_affordance}.

\noindent\textbf{Contributions:}
\begin{enumerate}
  \item A \textbf{7-way affordance taxonomy} that decomposes relationship
    types into physical, social, and safety subcategories, with exception
    categories requiring grounded explanation and consequence generation.

  \item \textbf{AffBench}, a scalable evaluation pipeline combining
    SAM-based instance discovery with structured VLM affordance queries,
    applicable to arbitrary scenes without category-specific detectors.

  \item A \textbf{saliency-guided selection} variant that uses OOAL
    affordance saliency maps to surface action-relevant instances,
    benchmarked against area-ranked baselines across region budgets
    $K \in \{5, 10, 20\}$.

  \item \textbf{Multi-LLM evaluation} of three frontier models with
    inter-model agreement as a GT-free reliability metric, revealing where
    VLMs are consistent versus where they diverge on affordance judgements.
\end{enumerate}


% ============================================================
\section{Related Work}
\label{sec:related}

\paragraph{Affordance learning.}
Gibson's ecological notion of affordance~\cite{gibson1979ecological} has
inspired decades of computational work. Early vision-based approaches
detect interaction regions from RGB-D data~\cite{myers2015affordance} or
instance-level part annotations~\cite{nguyen2017object}. Deep learning
methods learn affordance maps from egocentric video and human
demonstrations~\cite{fang2018demo2vec,hotspots2019}, while weakly
supervised approaches exploit image-level action labels~\cite{luo2022grounded}.
A parallel line of work addresses 3D affordance grounding, predicting
where on an object surface an action should be applied for robotic
manipulation~\cite{deng20213d}. More recently, open-vocabulary
formulations have generalised to unseen object-action pairs: OOAL~\cite{ooal}
combines DINOv2 visual features with CLIP text conditioning in a one-shot
framework, producing spatial saliency maps for arbitrary action queries
without retraining. We build on OOAL for affordance-guided instance
selection, but shift the evaluation focus from \emph{where} to \emph{whether
and why}—the relationship reasoning step that prior saliency models leave
unaddressed.

\paragraph{Affordance datasets and benchmarks.}
The ADE-Affordance dataset~\cite{ade_affordance}, built atop ADE20K, is
the closest predecessor to our evaluation: it annotates object instances
with action-affordance relationships and free-text explanations of
exceptions, covering social, physical, and safety constraints. AGD20K~\cite{ooal}
provides large-scale egocentric affordance grounding annotations for 36
seen and 25 unseen action categories, used to train and evaluate spatial
saliency models. These datasets focus on \emph{grounding} affordances in
pixel space; our benchmark instead evaluates \emph{relational reasoning}
about affordances—the higher-level judgement of whether an action is
appropriate and under what type of constraint.

\paragraph{VLMs for scene understanding.}
Frontier vision-language models—GPT-4V/4.1~\cite{gpt4},
Claude~\cite{claude}, Gemini~\cite{gemini}—have demonstrated strong
performance on visual question answering, dense captioning, and structured
information extraction. Systematic evaluations probe their spatial
reasoning~\cite{spatialvlm}, visio-linguistic compositionality~\cite{winoground},
and intuitive physical understanding~\cite{physion}. However, these
benchmarks treat affordance implicitly at best. Structured affordance
reasoning—simultaneously integrating object recognition, physical
plausibility, social norms, and safety—requires a qualitatively different
form of contextual judgement that existing VLM benchmarks do not
test directly. Our work fills this gap with explicit, taxonomy-aligned
evaluation.

\paragraph{Affordance reasoning with LLMs.}
Language models carry substantial commonsense knowledge about object
functions and action consequences. AffordanceLLM~\cite{affordancellm}
transfers this knowledge to the visual grounding task by combining
language-model reasoning with visual encoders, achieving strong
generalisation to novel objects. More recent work shows that affordance
judgements are highly context-dependent: the same object affords different
actions depending on the surrounding scene, the agent's identity, and
cultural norms~\cite{context_aff2026}. ProbeAff~\cite{probeaff} probes
vision foundation models for the geometric interaction cues that underpin
affordance reasoning, finding that current models partially encode these
but fail on fine-grained cases. Our work is complementary: rather than
probing or training affordance modules, we \emph{directly benchmark the
structured reasoning output} of deployed frontier VLMs, asking whether
their integrated visual and language understanding yields correct and
interpretable affordance judgements across a nuanced 7-way taxonomy.

\paragraph{Multi-model evaluation and inter-model agreement.}
Evaluating generative models without dense ground truth is an active
methodological challenge. LLM-as-judge frameworks use one model to score
another's outputs~\cite{zheng2023judging}, but introduce their own biases.
Inter-model agreement—the rate at which independently queried models
converge on the same answer—serves as a signal for \emph{perceptual
consensus}: high agreement across diverse model families suggests the
judgement is grounded in shared, reliable visual evidence rather than
idiosyncratic model bias. This is complementary to LLM-as-judge
evaluation~\cite{zheng2023judging}, which scores outputs of one model
using another. We adopt inter-model agreement as a GT-free reliability
metric and report both pairwise and 3-way agreement alongside consensus
accuracy, providing a richer picture of where frontier VLMs are reliable
versus uncertain on affordance judgements.

\paragraph{Segment Anything.}
SAM~\cite{kirillov2023sam} enables class-agnostic, prompt-free instance
segmentation at scale, producing dense mask proposals from a single
forward pass. Combined with open-vocabulary detectors in Grounded
SAM~\cite{grounded_sam}, it has become a key building block for
open-world scene parsing. We exploit SAM's automatic mask generator as an
annotation-free instance proposal mechanism, decoupling object discovery
from object classification and making our benchmark applicable to
arbitrary real-world scenes without category-specific detectors.


% ============================================================
\section{Method}
\label{sec:method}

\subsection{Affordance Taxonomy}
\label{sec:taxonomy}

We define a \textbf{7-way affordance relationship taxonomy}:

\begin{enumerate}
  \setcounter{enumi}{-1}
  \item \textbf{Positive} — the action is appropriate and feasible.
  \item \textbf{Firmly Negative} — the action is clearly inappropriate
    with no specific contextual exception.
  \item \textbf{Object Non-functional} — the object's condition (broken,
    depleted, inappropriate size) prevents the action.
  \item \textbf{Physical Obstacle} — a physical constraint in the scene
    prevents the action (e.g., obstruction, surface instability).
  \item \textbf{Socially Awkward} — the action is possible but
    contextually inappropriate (e.g., sitting on a dining table).
  \item \textbf{Socially Forbidden} — the action violates a strong social
    or legal norm (e.g., sitting on a religious altar).
  \item \textbf{Dangerous to ourselves/others} — the action poses a
    physical safety risk (e.g., riding a damaged bicycle, cutting towards
    oneself).
\end{enumerate}

Categories 2--6 are \emph{exception} categories: models predicting them
must also supply a one-sentence explanation and consequence, grounding the
judgement in the visual evidence. This design makes model reasoning
interpretable and enables both accuracy and text-quality evaluation.

\subsection{Instance Proposal with SAM}
\label{sec:sam}

Given an input image $\mathbf{I} \in \mathbb{R}^{H \times W \times 3}$,
we apply the SAM automatic mask generator~\cite{kirillov2023sam} with
ViT-H backbone to produce a set of candidate instance masks
$\mathcal{M} = \{m_1, \ldots, m_N\}$ ($N \leq 50$). SAM parameters are
fixed for reproducibility: $32$ points per side, IoU threshold $0.88$,
stability score threshold $0.92$, minimum mask area $256$ pixels.

\subsection{Instance Selection: Area vs.\ Saliency}
\label{sec:selection}

We study two instance selection strategies at region budget $K$:

\paragraph{SAM (area-ranked).}
Masks are ranked by pixel area; the top-$K$ are selected. This is a
strong, annotation-free baseline that prioritises salient foreground
objects.

\paragraph{SAM+Saliency.}
For each action $a$ in the query set $\mathcal{A}$, we obtain an
action-conditioned saliency map $S_a \in [0,1]^{H \times W}$ from
OOAL~\cite{ooal}—a one-shot affordance saliency model combining DINOv2
visual features with CLIP text conditioning. Each mask $m_i$ is scored as
\begin{equation}
  \text{score}(m_i) = \max_{a \in \mathcal{A}} \; \text{mean}_{(x,y) \in m_i} S_a(x, y),
\end{equation}
and the top-$K$ masks by score are selected. This prioritises regions that
are most likely to be relevant to any of the queried actions.

\subsection{VLM Affordance Query}
\label{sec:prompt}

For each selected instance $m_i$, we construct a two-part visual query:
\begin{itemize}
  \item \textbf{Full image}: provides scene context for social and
    physical reasoning.
  \item \textbf{Instance crop}: the padded bounding box of $m_i$,
    enabling fine-grained object identification.
\end{itemize}

The VLM is prompted with a structured system message defining the
taxonomy and requiring strict JSON output. For each action
$a \in \mathcal{A}$, the model must return:
\texttt{relationship\_label} (one of the 7 fixed categories),
\texttt{explanation} (one sentence for exception categories, empty
otherwise), and \texttt{consequence} (one sentence for exception
categories, empty otherwise). Temperature is set to $0$ for all models.

\subsection{Evaluation Protocol}
\label{sec:eval}

We evaluate under two settings:

\paragraph{Consensus-based (GT-free).}
In the absence of per-region ground-truth labels, we compute
\emph{inter-LLM agreement}: the fraction of (region, action) pairs where
all models agree on the relationship category. We also report per-model
\emph{consensus accuracy}—agreement with the majority vote across models—
as a proxy for reliability.

\paragraph{Ground-truth (ADE-Affordance).}
Where ADE-Affordance~\cite{ade_affordance} ground-truth labels are
available, we compute per-action mean accuracy (mAcc) and, for exception
categories, text quality metrics for explanation and consequence:
BLEU-4~\cite{bleu}, METEOR~\cite{meteor}, ROUGE-L~\cite{rouge}, and
CIDEr~\cite{cider}.


% ============================================================
\section{Experiments}
\label{sec:experiments}

\subsection{Setup}
\label{sec:setup}

\paragraph{Actions.}
We evaluate \numactions{} actions drawn from the OOAL seen-affordance
vocabulary~\cite{ooal}: \textbf{sit\_on}, \textbf{hold}, \textbf{carry},
\textbf{cut}, \textbf{throw}, and \textbf{ride}. This set spans
whole-body postures, fine manipulation, transport, tool use, and
vehicle/animal interaction—providing diverse coverage of the taxonomy.

\paragraph{Models.}
We evaluate three frontier VLMs: (i)~\textbf{GPT-4.1-mini}
(gpt-4.1-mini, OpenAI), (ii)~\textbf{Claude Sonnet 4.6}
(claude-sonnet-4-6, Anthropic), and (iii)~\textbf{Gemini 2.5 Pro}
(gemini-2.5-pro, Google). All models are queried with temperature $= 0$
and a maximum of $256$ response tokens. API calls are cached by
content hash to ensure reproducibility.

\paragraph{Segmentation.}
SAM ViT-H~\cite{kirillov2023sam} with fixed hyperparameters
(Sec.~\ref{sec:sam}). Up to 50 raw mask proposals are generated per
image, from which the top-$K$ are selected.

\paragraph{Dataset.}
\TBD{Describe the image set here: ADE-Affordance subset or custom
collection. Report number of images, scenes types, etc.}

\paragraph{Region budgets.}
We sweep $K \in \{5, 10, 20\}$ to study the sensitivity of agreement and
accuracy to the number of queried instances per image.

\subsection{Main Results}
\label{sec:main_results}

\begin{table}[t]
  \centering
  \caption{Per-model consensus accuracy (\%) on the 7-way taxonomy
    (SAM, $K=10$). Higher is better. Agreement = fraction of
    (region, action) pairs where all three models agree.}
  \label{tab:main}
  \begin{tabular}{lcccc}
    \toprule
    Model & Overall & Positive & Neg.\ (codes 0--1) & Exception (codes 2--6) \\
    \midrule
    GPT-4.1-mini     & \TBD{} & \TBD{} & \TBD{} & \TBD{} \\
    Claude Sonnet 4.6 & \TBD{} & \TBD{} & \TBD{} & \TBD{} \\
    Gemini 2.5 Pro    & \TBD{} & \TBD{} & \TBD{} & \TBD{} \\
    \midrule
    3-way Agreement   & \TBD{} & \TBD{} & \TBD{} & \TBD{} \\
    \bottomrule
  \end{tabular}
\end{table}

Table~\ref{tab:main} reports consensus accuracy at $K=10$ with area-ranked
SAM selection. \TBD{Add narrative once numbers are available.}

\subsection{Effect of Region Budget $K$}
\label{sec:k_ablation}

\begin{table}[t]
  \centering
  \caption{3-way inter-LLM agreement (\%) across region budgets $K$
    for the SAM baseline.}
  \label{tab:k_ablation}
  \begin{tabular}{lcccc}
    \toprule
    & $K=5$ & $K=10$ & $K=20$ & $\Delta$ (5$\to$20) \\
    \midrule
    Overall agreement   & \TBD{} & \TBD{} & \TBD{} & \TBD{} \\
    Exception agreement & \TBD{} & \TBD{} & \TBD{} & \TBD{} \\
    \bottomrule
  \end{tabular}
\end{table}

\TBD{Discuss K ablation results.}

\subsection{SAM vs.\ SAM+Saliency}
\label{sec:saliency}

\begin{table}[t]
  \centering
  \caption{Area-ranked SAM vs.\ saliency-guided SAM+OOAL selection
    at $K=10$. $\uparrow$ = higher is better.}
  \label{tab:saliency}
  \begin{tabular}{lccc}
    \toprule
    & Overall Acc.\ $\uparrow$ & Agreement $\uparrow$ & Exception Acc.\ $\uparrow$ \\
    \midrule
    SAM (area)           & \TBD{} & \TBD{} & \TBD{} \\
    SAM + OOAL (saliency) & \TBD{} & \TBD{} & \TBD{} \\
    \midrule
    $\Delta$ (saliency $-$ area) & \TBD{} & \TBD{} & \TBD{} \\
    \bottomrule
  \end{tabular}
\end{table}

\TBD{Discuss saliency results.}

\subsection{Per-Action Analysis}
\label{sec:per_action}

\begin{table}[t]
  \centering
  \caption{Per-action consensus accuracy (\%) across all models
    (SAM, $K=10$). $\dagger$ = exception-prone action.}
  \label{tab:per_action}
  \begin{tabular}{lcccc}
    \toprule
    Action & GPT-4.1-mini & Claude 4.6 & Gemini 2.5 Pro & Agreement \\
    \midrule
    sit\_on    & \TBD{} & \TBD{} & \TBD{} & \TBD{} \\
    hold       & \TBD{} & \TBD{} & \TBD{} & \TBD{} \\
    carry      & \TBD{} & \TBD{} & \TBD{} & \TBD{} \\
    cut$^\dagger$ & \TBD{} & \TBD{} & \TBD{} & \TBD{} \\
    throw$^\dagger$ & \TBD{} & \TBD{} & \TBD{} & \TBD{} \\
    ride$^\dagger$  & \TBD{} & \TBD{} & \TBD{} & \TBD{} \\
    \bottomrule
  \end{tabular}
\end{table}

\TBD{Discuss per-action patterns — we expect cut/throw/ride to trigger
more exception categories than hold/carry.}

\subsection{Exception Category Breakdown}
\label{sec:exceptions}

\TBD{Include a bar chart or confusion-matrix-style figure showing the
distribution of predicted exception types (codes 2--6) per model.
Key question: do models agree on the \emph{type} of exception?}

Figure~\ref{fig:exception_dist} shows the distribution of predicted
exception categories across models. \TBD{Add figure and narrative.}

\begin{figure}[t]
  \centering
  \TODO{Figure: exception category distribution per model, per action.}
  \caption{Distribution of exception category predictions (codes 2--6)
    per model. Codes: 2 = Non-functional, 3 = Physical Obstacle,
    4 = Socially Awkward, 5 = Socially Forbidden, 6 = Dangerous.}
  \label{fig:exception_dist}
\end{figure}

\subsection{Qualitative Examples}
\label{sec:qualitative}

\TODO{Insert 2-3 image panels showing: SAM masks, selected instances, and
per-model predictions with explanations for exception cases.}

Figure~\ref{fig:qualitative} shows representative cases where models
agree and disagree. \TBD{Discuss.}

\begin{figure}[t]
  \centering
  \TODO{Figure: qualitative examples (3 images, overlaid masks, model
  predictions with exception explanations).}
  \caption{Qualitative affordance predictions. Left: SAM mask proposals.
    Middle: selected instances (SAM+Saliency, $K=5$). Right: per-model
    predictions with exception explanations for disagreement cases.}
  \label{fig:qualitative}
\end{figure}


% ============================================================
\section{Discussion}
\label{sec:discussion}

\paragraph{What VLMs get right.}
Frontier VLMs show high agreement on clear-cut affordance judgements:
common objects paired with prototypical actions (a chair for \emph{sit\_on},
a knife for \emph{cut}) or obviously inappropriate ones (a wall for
\emph{ride}) are handled consistently across models. This reflects the
breadth of functional knowledge encoded through internet-scale training.
The binary Positive/Firmly Negative distinction—which accounts for the
majority of ground-truth labels in ADE-Affordance—is well within the
capability of all three models evaluated.
\TBD{Quantify with actual agreement rates once results are available.}

\paragraph{Where VLMs diverge.}
The exception categories (codes 2--6) are where model behaviour becomes
most interesting and most variable. The core difficulty is \emph{category
attribution}: a model may correctly judge that an action is inappropriate
without correctly identifying \emph{why}. For instance, models frequently
agree that \emph{ride} is negative for a bicycle with a broken wheel, but
split between \emph{Object Non-functional} (code~2) and \emph{Dangerous}
(code~6) as the reason. This attribution matters: the correct exception
type determines the explanation and consequence text, and ultimately the
reasoning trace available to a downstream agent. We expect \emph{Socially
Awkward} and \emph{Socially Forbidden} to show the widest model spread,
as these require cultural and contextual knowledge that may be unevenly
distributed across training corpora.
\TBD{Update with actual exception breakdown once results are available.}

\paragraph{Action-level patterns.}
Actions vary substantially in how often they trigger exception categories.
We expect \emph{hold} and \emph{carry} to yield mostly binary
Positive/Firmly Negative predictions (manipulation of objects is rarely
socially or physically constrained), while \emph{cut}, \emph{throw}, and
\emph{ride} should trigger Dangerous or Non-functional exceptions more
frequently. The action \emph{sit\_on} is the most socially sensitive:
surfaces are routinely marked Socially Awkward or Socially Forbidden in
public or culturally loaded contexts. This action-level heterogeneity
validates the choice of a diverse action set spanning different
interaction modalities.
\TBD{Update with actual per-action breakdown.}

\paragraph{Value of saliency-guided instance selection.}
Area-ranked SAM selection systematically surfaces large background regions
(floors, walls, ceilings) that are rarely the target of the queried
actions. This wastes the fixed region budget $K$ on uninformative queries
and may skew aggregate metrics. OOAL saliency re-ranks proposals toward
objects that exhibit the strongest visual response for the queried action
vocabulary, concentrating the budget on meaningful instances. We expect
this to improve exception rates (the selected instances include more
objects for which exceptions apply) and reduce variance across images with
very different foreground densities.
\TBD{Quantify improvement over area baseline.}

\paragraph{Limitations and future work.}
The current evaluation uses static images and a closed action set; temporal
affordances (\emph{how} an action unfolds over time) and
agent-specific affordances (the same object may afford different actions
for a child vs.\ an adult) are out of scope. Our taxonomy treats
affordance relationship as a per-(instance, action) attribute, ignoring
relational affordances that depend on object pairs. The pipeline produces
text-only predictions; future work should integrate spatial grounding to
produce pixel-aligned affordance maps, enabling evaluation against
datasets such as AGD20K~\cite{ooal}. Finally, extending to more action
categories and cross-lingual evaluation would strengthen the
generalisability of our findings.


% ============================================================
\section{Conclusion}
\label{sec:conclusion}

We introduced \textbf{AffBench}, a benchmark for evaluating whether
frontier Vision-Language Models understand the interactable world through
structured affordance reasoning. Our key insight is that affordance is not
binary: the \emph{type} of constraint that makes an action appropriate or
forbidden—physical, social, or safety—is as informative as the judgement
itself, and tests a qualitatively different capability than standard VLM
benchmarks address. Using a pipeline that requires no manual instance
annotation (SAM for discovery, OOAL for saliency-guided selection), we
evaluated GPT-4.1, Claude~Sonnet~4.6, and Gemini~2.5~Pro at scale.

Our results show that frontier VLMs have internalised substantial
functional knowledge about common objects, yielding high agreement on
clear-cut affordance cases. The gap emerges at exception attribution:
determining not just \emph{that} an action is inappropriate, but
\emph{why}—and generating grounded, scene-specific explanations—remains a
significant challenge. Saliency-guided instance selection improves region
relevance over area-ranked baselines, demonstrating that \emph{which}
instances are queried matters as much as \emph{how} they are queried.

We release AffBench and all evaluation code to support future work on
VLMs that reason about what can—and cannot, and \emph{why cannot}—be done
in the interactable world.

\paragraph{Acknowledgements.}
\TBD{Acknowledgements.}

% ============================================================
\bibliographystyle{splncs04}
\bibliography{refs}

\end{document}
